% SPDX-License-Identifier: PMPL-1.0-or-later
% SPDX-FileCopyrightText: 2026 Jonathan D.A. Jewell
%
% Typing the Boundaries: Preserving Type Safety Across
% Every System Crossing Point
%
% Intended venue: arXiv preprint (cs.PL / cs.SE / cs.NI)
% Subsequently: POPL, OOPSLA, or NSDI
%
% This is a COMPANION paper to the Groove Protocol paper.
% Groove is one mechanism; this paper presents the general thesis.

\documentclass[11pt,a4paper]{article}
\usepackage[utf8]{inputenc}
\usepackage[T1]{fontenc}
\usepackage{amsmath,amssymb,amsthm}
\usepackage{listings}
\usepackage{hyperref}
\usepackage{booktabs}
\usepackage{graphicx}
\usepackage{xcolor}
\usepackage[margin=2.5cm]{geometry}

\definecolor{codegreen}{rgb}{0.25,0.49,0.31}
\definecolor{codegray}{rgb}{0.5,0.5,0.5}
\definecolor{codeblue}{rgb}{0.22,0.40,0.72}

\lstset{
  basicstyle=\ttfamily\small,
  keywordstyle=\color{codeblue}\bfseries,
  commentstyle=\color{codegray}\itshape,
  stringstyle=\color{codegreen},
  breaklines=true,
  frame=single,
  captionpos=b
}

\newtheorem{theorem}{Theorem}
\newtheorem{lemma}{Lemma}
\newtheorem{definition}{Definition}
\newtheorem{property}{Property}
\newtheorem{observation}{Observation}

\title{Typing the Boundaries:\\
Preserving Type Safety Across Every System Crossing Point}

\author{Jonathan D.A. Jewell\\
\texttt{j.d.a.jewell@open.ac.uk}\\
The Open University}

\date{March 2026}

\begin{document}

\maketitle

\begin{abstract}
Modern type systems provide strong guarantees within a single language,
runtime, or process. But the moment data crosses a boundary --- between
services, between protocol layers, between browser and native code,
between network segments --- types evaporate. Data becomes untyped bytes,
and correctness reverts to runtime parsing, hope, and prayer.

We present an architecture for \emph{preserving type safety across every
system crossing point}. Using dependent types (Idris2) for interface
definitions, zero-overhead FFI (Zig) for crossing the language boundary,
structural capability matching (Groove Protocol) for crossing the network
boundary, frame-level routing with formal proofs (Typed Frame Router) for
crossing the protocol boundary, and clade-based provisioning (FLI) for
crossing the component boundary, we demonstrate that types can survive
every crossing they encounter --- from kernel socket to browser DOM.

We identify six boundary categories where types traditionally die,
present a mechanism for each that preserves type safety through the
crossing, and prove that the composed system of boundary-crossing
mechanisms is itself type-safe. The result is a world where the gaps
between systems are typed, not just the systems themselves.
\end{abstract}

\section{Introduction}

\subsection{The Boundary Problem}

Every mature programming language provides type safety within its own
domain. Rust guarantees memory safety. Haskell guarantees referential
transparency. Idris2 guarantees totality and resource linearity. Within
a single program, written in a single language, running in a single
process, types work.

But software does not live in a single process. It lives in the gaps
between processes, between machines, between networks, between runtimes.
And at every gap, types die.

\begin{observation}[The Type Death Problem]
When data crosses a system boundary, static type information is lost.
The receiving system must reconstruct types from untyped bytes via
parsing, validation, and runtime checks --- all of which can fail.
\end{observation}

Consider a simple example: a Rust service sends a JSON response to a
TypeScript frontend. Rust's type system guarantees the response is
well-formed at compile time. But the moment the response crosses the
HTTP boundary, it becomes an untyped byte stream. The TypeScript
frontend parses it, casts it to an interface, and hopes the shape
matches. If it doesn't, the error manifests at runtime --- in
production, at 3am.

This is not a Rust problem or a TypeScript problem. It is a
\emph{boundary problem}. Every boundary between systems is a gap where
types die and must be resurrected, imperfectly, on the other side.

\subsection{The Thesis}

We argue that this is not inevitable. Types can be made to survive
every crossing point, if the boundary itself is typed.

\begin{definition}[Typed Boundary]
A system boundary is \emph{typed} if there exists a compile-time proof
that data passing through the boundary preserves its structural type.
The proof is checked before the boundary is crossed, not after.
\end{definition}

We present six boundary categories and a mechanism for typing each one:

\begin{table}[h]
\centering
\small
\begin{tabular}{llll}
\toprule
Boundary & Crossing & Mechanism & Proof Technology \\
\midrule
Language & Rust $\leftrightarrow$ Zig $\leftrightarrow$ C
  & ABI/FFI triple & Idris2 dependent types \\
Service & Service A $\leftrightarrow$ Service B
  & Groove Protocol & Structural subtyping \\
Protocol & IPv4 $\leftrightarrow$ IPv6
  & Typed Frame Router & Transport transparency \\
Component & Panel A $\leftrightarrow$ FLI module
  & Clade inheritance & Trait matching \\
Runtime & Browser $\leftrightarrow$ Native
  & Groove Harness & CORS + type proofs \\
Storage & App $\leftrightarrow$ Database
  & VeriSimDB octads & 8-modality typing \\
\bottomrule
\end{tabular}
\caption{Six typed boundaries}
\label{tab:boundaries}
\end{table}

\subsection{Contributions}

\begin{enumerate}
\item A taxonomy of six boundary categories where types traditionally
  die, with a concrete mechanism for typing each (Section~\ref{sec:taxonomy}).
\item A three-layer architecture (ABI/FFI/API) that preserves dependent
  type proofs across the language boundary with zero runtime overhead
  (Section~\ref{sec:language}).
\item The Groove Protocol for preserving type safety across the service
  boundary via structural capability matching
  (Section~\ref{sec:service}).
\item The Typed Frame Router for preserving type safety across the
  protocol boundary via formally verified frame translation
  (Section~\ref{sec:protocol}).
\item FLI (Frontend-user Library Interface) for preserving type safety
  across the component boundary via clade-based capability inheritance
  (Section~\ref{sec:component}).
\item A composition theorem proving that the combined system of
  boundary-crossing mechanisms is itself type-safe --- types that enter
  at one boundary exit intact at another (Section~\ref{sec:composition}).
\end{enumerate}

\section{The Boundary Taxonomy}
\label{sec:taxonomy}

We identify six categories of boundary where types die in practice.
Each represents a different kind of crossing, with different failure
modes and different requirements for the typing mechanism.

\subsection{Boundary 1: Language}

\textbf{Crossing:} Data moves between code written in different
programming languages (e.g., Idris2 to Zig to C to Rust).

\textbf{How types die:} Each language has its own type system with
different representations, alignment rules, and calling conventions.
A Haskell algebraic data type has no direct representation in C.

\textbf{Traditional solution:} FFI with manual marshalling. Developers
write conversion code by hand. Errors are common and caught at runtime.

\textbf{Our mechanism:} The ABI/FFI/API triple.
\begin{itemize}
\item \textbf{Idris2 ABI} defines the canonical interface with dependent
  types. Every type has a tag encoding with a roundtrip proof: encoding
  then decoding is the identity (proven at compile time).
\item \textbf{Zig FFI} implements the C-compatible functions. Zig's
  explicit memory layout guarantees match the Idris2 specifications.
\item \textbf{V API} generates ergonomic bindings for target languages,
  carrying the type constraints to the surface.
\end{itemize}

Types survive because the \emph{boundary itself} (the C ABI) is typed
by Idris2. The proof that the encoding is lossless is checked before
the code compiles.

\subsection{Boundary 2: Service}

\textbf{Crossing:} Data moves between independent services over the
network (e.g., a desktop app and a voice platform).

\textbf{How types die:} Services communicate via HTTP/JSON/gRPC. The
type information exists in protocol buffer definitions or OpenAPI specs,
but compliance is checked at runtime, not compile time.

\textbf{Traditional solution:} API contracts (OpenAPI, protobuf).
Schema validation at runtime. Contract testing in CI.

\textbf{Our mechanism:} The Groove Protocol. Services declare typed
capability interfaces in Idris2. Structural subtyping determines
compatibility. If the types don't match, the connection \emph{cannot be
expressed} in the type system --- it is a compile-time impossibility,
not a runtime error.

\subsection{Boundary 3: Protocol}

\textbf{Crossing:} Data moves between different network protocols
(e.g., IPv4 to IPv6, Ethernet to Fibre Channel).

\textbf{How types die:} Protocol translation is traditionally done by
untyped packet manipulation. Routers, NAT boxes, and protocol converters
operate on raw bytes with no formal guarantee that the translation
preserves the data.

\textbf{Traditional solution:} NAT64, protocol gateways, manual
testing.

\textbf{Our mechanism:} The Typed Frame Router. A formally verified
frame-level proxy that proves four properties:
\begin{enumerate}
\item \emph{Transport transparency}: bytes in = bytes out
\item \emph{Bounded memory}: $\leq$ 4,096 bytes per connection
\item \emph{Direction safety}: translation direction cannot be reversed
\item \emph{Connection liveness}: every connection terminates
\end{enumerate}

Implemented in Zig with zero-copy \texttt{splice(2)} on Linux. The
proofs are in Idris2; the runtime overhead is zero.

\subsection{Boundary 4: Component}

\textbf{Crossing:} Data moves between UI components within an
application (e.g., a panel requesting a gauge library).

\textbf{How types die:} Components typically communicate via events,
callbacks, or shared state. The interface between components is
implicit --- a function signature, at best. Most component systems
(React, Vue, Web Components) have no formal mechanism for proving
that a component's required interface is satisfied by its provider.

\textbf{Traditional solution:} PropTypes, TypeScript interfaces,
runtime checks.

\textbf{Our mechanism:} FLI (Frontend-user Library Interface) with
clade-based inheritance. Components declare their required capabilities
as clade traits (e.g., \texttt{fli-gauge}, \texttt{fli-editable}).
The provisioner matches traits to modules. If a module is missing,
the component falls back to base functionality --- no crash, no error.
Capability fit is structural, following the Groove pattern.

\subsection{Boundary 5: Runtime}

\textbf{Crossing:} Data moves between the browser security sandbox
and native desktop code.

\textbf{How types die:} Browser extensions communicate with local
services via \texttt{fetch()} to localhost. The request is untyped
HTTP. The response is parsed JSON. Neither side has compile-time proof
that the other conforms to the expected interface.

\textbf{Traditional solution:} Chrome Native Messaging, manual JSON
validation, hope.

\textbf{Our mechanism:} The Groove Browser Harness. A browser extension
that probes localhost for Groove manifests and establishes type-safe
connections. The Idris2 ABI types cross the browser/native boundary via
the C header $\rightarrow$ Zig FFI $\rightarrow$ Groove HTTP
$\rightarrow$ \texttt{fetch()} chain. At each hop, the types are
preserved.

\subsection{Boundary 6: Storage}

\textbf{Crossing:} Data moves between application code and a database.

\textbf{How types die:} ORMs map objects to tables. The mapping is
implicit and fragile. Schema migrations can silently break the mapping.
NoSQL databases have no schema at all.

\textbf{Traditional solution:} ORMs, schema migrations, runtime
validation.

\textbf{Our mechanism:} VeriSimDB octads. Every stored record has eight
typed modalities (document, semantic, temporal, spatial, graph, vector,
provenance, binary). The type of each modality is declared in the
schema. Drift detection continuously verifies that stored data matches
its declared type. Divergence triggers alerts, not silent corruption.

\section{Preserving Types Across the Language Boundary}
\label{sec:language}

The language boundary is the most fundamental. If types cannot survive
the crossing from one language to another, no higher-level boundary
mechanism can be trusted.

Our approach uses three layers:

\begin{enumerate}
\item \textbf{Idris2 ABI}: defines the interface with dependent types.
  Every sum type gets a tag encoding function and a decoding function,
  connected by a \emph{roundtrip proof}:

\begin{lstlisting}[language=Haskell]
-- Encoding
frameFamilyToTag : FrameFamily -> Bits8
frameFamilyToTag IPv4 = 0
frameFamilyToTag IPv6 = 1

-- Decoding
tagToFrameFamily : Bits8 -> Maybe FrameFamily
tagToFrameFamily 0 = Just IPv4
tagToFrameFamily 1 = Just IPv6
tagToFrameFamily _ = Nothing

-- Proof: decode(encode(x)) = Just x
frameFamilyRoundtrip : (f : FrameFamily)
    -> tagToFrameFamily (frameFamilyToTag f) = Just f
frameFamilyRoundtrip IPv4 = Refl
frameFamilyRoundtrip IPv6 = Refl
\end{lstlisting}

  The \texttt{Refl} proof term is erased at runtime. The guarantee is
  free.

\item \textbf{Zig FFI}: implements the C-compatible functions. Zig's
  explicit control over memory layout means the runtime representation
  exactly matches the Idris2 specification. No hidden padding, no
  alignment surprises, no garbage collector interference.

\item \textbf{V API}: generates ergonomic bindings for each target
  language. The bindings enforce the type constraints at the surface
  level, even in languages that cannot express dependent types.
\end{enumerate}

\section{Preserving Types Across the Service Boundary}
\label{sec:service}

See companion paper: \emph{Groove: Formally Verified System Composition
via Dependent Types and Linear Resources} \cite{groove2026}.

The key insight is that service compatibility is expressed as structural
subtyping. The \texttt{auto} proof search in Idris2 automatically finds
(or fails to find) a proof that one service's offered capabilities
satisfy another's requirements. Failure is a compile-time error, not a
runtime exception.

\section{Preserving Types Across the Protocol Boundary}
\label{sec:protocol}

The Typed Frame Router sits at any network boundary and translates
frames between protocol families while preserving four formally verified
properties.

The critical proof is \emph{transport transparency}: for any byte
sequence $b$ entering the router, the same byte sequence $b$ exits.
The router is a pure conduit. Combined with the ABI roundtrip proofs,
this means typed data that enters as Idris2 values, crosses the C ABI,
crosses the network via IPv4, passes through the Typed Frame Router to
IPv6, and arrives at the destination --- is \emph{provably identical}
to what was sent.

The types survived the protocol boundary because the boundary itself
is typed.

\section{Preserving Types Across the Component Boundary}
\label{sec:component}

FLI follows the Groove pattern at the component level: panels declare
traits, the provisioner matches modules to traits, and unmatched traits
result in graceful fallback rather than errors. The key property is:

\begin{property}[Component Bare-First]
Every panel functions correctly with zero FLI modules loaded. FLI
capabilities are additive. No module is required.
\end{property}

This is the same property as Groove's bare-first principle, applied
to the component boundary rather than the service boundary.

\section{Composition: Types Across Multiple Boundaries}
\label{sec:composition}

The most important result is that these mechanisms compose. A value
can cross multiple boundaries and retain its type throughout:

\begin{lstlisting}[caption={A value crossing four boundaries}]
Idris2 value
  -> [ABI roundtrip proof] -> Zig C function
  -> [Groove structural match] -> HTTP to remote service
  -> [Typed Frame Router transparency] -> IPv6 from IPv4
  -> [FLI trait match] -> Browser panel via Groove Harness
\end{lstlisting}

\begin{theorem}[Boundary Composition]
If boundary $B_1$ preserves type $T$ (proven by $p_1$) and boundary
$B_2$ preserves type $T$ (proven by $p_2$), then the composed crossing
$B_1 \circ B_2$ preserves type $T$ (proven by transitivity of equality).
\end{theorem}

\begin{proof}
Each boundary mechanism produces a proof of the form $\text{output} =
\text{input}$. By transitivity, chaining $n$ such proofs yields
$\text{final\_output} = \text{original\_input}$. The proof terms are
erased at runtime; only the values flow.
\end{proof}

This means the entire chain --- from Idris2 value to browser DOM ---
is covered by a single composed proof. The types never die.

\section{Evaluation}

\subsection{Boundary Coverage}

\begin{table}[h]
\centering
\small
\begin{tabular}{lccc}
\toprule
Boundary & Types Survive? & Proof Mechanism & Runtime Cost \\
\midrule
Language (ABI/FFI) & \checkmark & Roundtrip lemma & Zero (erased) \\
Service (Groove) & \checkmark & Structural subtyping & Zero (compile-time) \\
Protocol (TFR) & \checkmark & Transparency proof & Zero (splice) \\
Component (FLI) & \checkmark & Trait matching & Near-zero (JS check) \\
Runtime (Browser) & \checkmark & CORS + type chain & Near-zero (fetch) \\
Storage (VeriSimDB) & \checkmark & Drift detection & Background (15s cycle) \\
\bottomrule
\end{tabular}
\caption{Type preservation across all six boundaries}
\end{table}

In five of six cases, the runtime cost of type preservation is zero or
near-zero. The proofs are compile-time artifacts that are erased before
execution. Only the storage boundary (VeriSimDB drift detection) has a
measurable ongoing cost, and it runs as a background process.

\subsection{What Cannot Be Typed}

We identify three boundaries where our approach does not apply:

\begin{enumerate}
\item \textbf{WAN/Internet}: we do not control the remote endpoint.
  Types can be preserved up to the network edge, but the public internet
  is untyped. The Typed Frame Router can operate at the egress point,
  but the far side may not participate.
\item \textbf{Encrypted tunnels (mid-stream)}: encrypted traffic cannot
  be inspected without terminating the encryption. The TFR can operate
  at the encryption boundary but not within the tunnel.
\item \textbf{Broadcast protocols}: connectionless broadcast has no
  point-to-point channel to type. A different primitive (typed frame
  \emph{filter}) would be needed.
\end{enumerate}

\section{Related Work}

\subsection{Gradual Typing}
Gradual typing \cite{siek2006} preserves types across the boundary
between typed and untyped code within a single language. Our work
extends this idea across \emph{system} boundaries, not just language
boundaries.

\subsection{Contract Systems}
Racket's contract system \cite{findler2002} enforces behavioural
contracts at module boundaries. Our approach is similar in spirit but
uses compile-time proofs rather than runtime monitors.

\subsection{Session Types}
Session types \cite{honda1993} type the communication \emph{protocol}
between two parties. Our work types the \emph{boundary} between
arbitrary systems, regardless of communication protocol.

\subsection{Dependent Interoperability}
Osera and Zdancewic \cite{osera2012} study dependent interoperability
between languages with different type systems. Our ABI/FFI/API triple
is a practical realisation of this idea, using Idris2 as the canonical
type language and C ABI as the universal crossing point.

\section{Discussion}

\subsection{Why Now?}

Three enabling technologies make boundary typing practical today:

\begin{enumerate}
\item \textbf{Idris2's quantitative type theory}: linear resource
  tracking at the type level, with efficient compilation and proof
  erasure.
\item \textbf{Zig's zero-overhead FFI}: a systems language that
  can match C ABI layout exactly, with no runtime, no GC, and
  comptime evaluation for embedding proofs.
\item \textbf{Widespread IPv6}: link-local addressing enables
  zero-configuration network discovery, making the Groove protocol
  practical without infrastructure.
\end{enumerate}

None of these existed in mature form five years ago. The convergence
makes boundary typing feasible for the first time.

\subsection{The World Is Not Untyped}

The common framing is that the world outside our programs is untyped ---
that networks carry raw bytes, databases store untyped blobs, and
interprocess communication is inherently unsafe.

We argue this framing is wrong. The world is \emph{typed at the edges
and untyped in the gaps between}. Each system has internal types. The
problem is not that types don't exist --- it is that they don't survive
the crossing from one system to another.

The solution is not to type the world. It is to type the boundaries.

\section{Conclusion}

We have presented an architecture for preserving type safety across six
categories of system boundary: language, service, protocol, component,
runtime, and storage. Each boundary has a specific mechanism ---
ABI roundtrip proofs, Groove structural matching, Typed Frame Router
transparency proofs, FLI clade provisioning, browser harness type
chains, and VeriSimDB drift detection --- and these mechanisms compose
via transitivity of equality proofs.

The result is that a typed value can flow from an Idris2 definition,
through a Zig FFI call, across a Groove connection, through a protocol
translation, into a browser extension, and arrive at the DOM ---
provably intact. The types never died. They just needed boundaries
that would carry them.

The world is not untyped. It is typed at the edges and untyped in the
gaps. We are closing the gaps.

\bibliographystyle{plain}
\begin{thebibliography}{10}

\bibitem{groove2026}
J.~D.~A.~Jewell.
\newblock Groove: Formally Verified System Composition via Dependent Types and
  Linear Resources.
\newblock arXiv preprint, 2026.

\bibitem{brady2021}
E.~Brady.
\newblock \emph{Idris 2: Quantitative Type Theory in Practice}.
\newblock In ECOOP, 2021.

\bibitem{siek2006}
J.~Siek and W.~Taha.
\newblock Gradual typing for functional languages.
\newblock In Scheme Workshop, 2006.

\bibitem{findler2002}
R.~Findler and M.~Felleisen.
\newblock Contracts for higher-order functions.
\newblock In ICFP, 2002.

\bibitem{honda1993}
K.~Honda.
\newblock Types for dyadic interaction.
\newblock In CONCUR, 1993.

\bibitem{osera2012}
P.~Osera and S.~Zdancewic.
\newblock Dependent interoperability.
\newblock In TLDI, 2012.

\end{thebibliography}

\end{document}
